\chapter{Popups and Dialogs}
\label{ch:popups}

Kiss provides a popup system for modal dialogs and several built-in dialog types for messages and confirmations.  Popups are the primary mechanism for editing records selected from a grid.

\section{The Popup Component}
\label{sec:popup-component}

A popup is defined in the screen's HTML using the \texttt{<popup>}\index{popup} tag.  It requires exactly two direct children: \texttt{<popup-title>} and \texttt{<popup-body>}.

\begin{lstlisting}[style=html]
<popup id="edit-popup" height="200px" width="450px">
    <popup-title>Edit Employee</popup-title>
    <popup-body>
        <div style="display: grid;
                    grid-template-columns: auto 1fr;
                    gap: 8px; padding: 10px;">
            <span>First Name:</span>
            <text-input id="ep-first" required
                        style="width: 250px;"></text-input>
            <span>Last Name:</span>
            <text-input id="ep-last" required
                        style="width: 250px;"></text-input>
        </div>
        <div style="text-align: right; margin-top: 10px;">
            <push-button id="ep-ok">OK</push-button>
            <push-button id="ep-cancel">Cancel</push-button>
        </div>
    </popup-body>
</popup>
\end{lstlisting}

\begin{warnbox}
The popup \emph{must} have exactly two direct children: \texttt{<popup-title>} and \texttt{<popup-body>}.  Using a single wrapper \texttt{<div>} instead causes a ``Cannot set properties of undefined'' error at runtime.

\begin{lstlisting}[style=html]
<!-- WRONG: causes error -->
<popup id="bad-popup">
    <div class="content">...</div>
</popup>

<!-- CORRECT -->
<popup id="good-popup">
    <popup-title>Title</popup-title>
    <popup-body>Content</popup-body>
</popup>
\end{lstlisting}
\end{warnbox}

\section{Opening and Closing Popups}
\label{sec:popup-openclose}

\begin{lstlisting}[style=javascript]
// Open popup, set focus to a control
Utils.popup_open('edit-popup', 'ep-first');

// Close popup
Utils.popup_close();
\end{lstlisting}

\noindent
\texttt{Utils.popup\_open()}\index{Utils.popup\_open} accepts two parameters:

\begin{enumerate}
  \item \textbf{Popup ID} --- the \texttt{id} of the \texttt{<popup>} element.
  \item \textbf{Focus control} (optional) --- the \texttt{id} of a component that should receive focus when the popup opens.
\end{enumerate}

\noindent
\texttt{Utils.popup\_close()}\index{Utils.popup\_close} closes the most recently opened popup.  You can optionally pass the popup ID for verification.

\subsection{Popup Stacking}

Multiple popups can be open simultaneously.  The framework manages a stack of open popups and increments the z-index for each new popup.  Closing a popup pops the stack and restores the previous popup's state.

\subsection{DOM Preservation After Close}

When a popup is closed, its DOM elements are \emph{hidden}, not destroyed.  This means:

\begin{itemize}
  \item All controls within the popup remain accessible.
  \item Values can be read from popup controls after closing.
  \item Control references (\texttt{\$\$('popup-control-id')}) continue to work.
  \item The popup can be reopened without re-initializing its controls.
\end{itemize}

\begin{lstlisting}[style=javascript]
Utils.popup_open('edit-popup');
// ... user interacts with popup ...
Utils.popup_close();

// Controls are still accessible after close
let name = $$('ep-first').getValue();  // works!
\end{lstlisting}

\section{Popup Sizing}
\label{sec:popup-sizing}

Specify both \texttt{height} and \texttt{width} attributes on the \texttt{<popup>} tag:

\begin{lstlisting}[style=html]
<popup id="my-popup" height="300px" width="500px">
\end{lstlisting}

\noindent
Size the popup appropriately for its content.  Avoid excess whitespace, and test on different screen sizes.

Dynamic sizing is also available:

\begin{lstlisting}[style=javascript]
Utils.popup_set_width('my-popup', '600px');
Utils.popup_set_height('my-popup', '400px');
\end{lstlisting}

\section{The Standard Edit Pattern}
\label{sec:popup-edit-pattern}

The most common popup use case is editing a record selected from a grid.  The pattern uses one popup for both adding and editing:

\begin{lstlisting}[style=javascript]
// New record
$$('new-btn').onclick(() => {
    $$('ep-title').setValue('Add Employee');
    $$('ep-first').clear();
    $$('ep-last').clear();
    Utils.popup_open('edit-popup', 'ep-first');

    $$('ep-ok').onclick(async () => {
        if ($$('ep-first').isError('First Name')) return;
        if ($$('ep-last').isError('Last Name')) return;
        const res = await Server.call(WS, 'create', {
            firstName: $$('ep-first').getValue(),
            lastName: $$('ep-last').getValue()
        });
        if (res._Success) {
            Utils.popup_close();
            loadGrid();
        }
    });

    $$('ep-cancel').onclick(() => Utils.popup_close());
});

// Edit selected record
$$('edit-btn').onclick(() => {
    const row = grid.getSelectedRow();
    $$('ep-title').setValue('Edit Employee');
    $$('ep-first').setValue(row.firstName);
    $$('ep-last').setValue(row.lastName);
    Utils.popup_open('edit-popup', 'ep-first');

    $$('ep-ok').onclick(async () => {
        if ($$('ep-first').isError('First Name')) return;
        if ($$('ep-last').isError('Last Name')) return;
        const res = await Server.call(WS, 'save', {
            id: row.id,
            firstName: $$('ep-first').getValue(),
            lastName: $$('ep-last').getValue()
        });
        if (res._Success) {
            Utils.popup_close();
            loadGrid();
        }
    });

    $$('ep-cancel').onclick(() => Utils.popup_close());
});
\end{lstlisting}

\begin{tipbox}
Use a \texttt{<text-label>} for the popup title so it can be changed programmatically between ``Add'' and ``Edit'' modes.
\end{tipbox}

\section{Message Dialogs}
\label{sec:popup-message}

\texttt{Utils.showMessage()}\index{Utils.showMessage} displays an informational dialog with an OK button:

\begin{lstlisting}[style=javascript]
Utils.showMessage('Info', 'Record saved successfully.');
Utils.showMessage('Error', 'Failed to connect to server.');
\end{lstlisting}

\noindent
The dialog is modal and draggable.  The title bar is colored based on the title text: red for ``Error'', blue otherwise.  The method returns a Promise that resolves when the user clicks OK.

\begin{lstlisting}[style=javascript]
await Utils.showMessage('Info', 'Please review the changes.');
// execution continues after user clicks OK
\end{lstlisting}

\section{Confirmation Dialogs}
\label{sec:popup-yesno}

\texttt{Utils.yesNo()}\index{Utils.yesNo} displays a Yes/No confirmation dialog:

\begin{lstlisting}[style=javascript]
// Callback style
Utils.yesNo('Confirm', 'Delete this record?',
    function () {
        // user clicked Yes
        deleteRecord();
    },
    function () {
        // user clicked No (optional)
    }
);

// Promise style
const answer = await Utils.yesNo('Confirm', 'Delete?');
if (answer) {
    deleteRecord();
}
\end{lstlisting}

\noindent
The dialog is modal, draggable, and mobile-responsive (adjusts width for small screens).

\section{Wait Messages}
\label{sec:popup-wait}

For long-running operations, display a non-modal wait message:

\begin{lstlisting}[style=javascript]
Utils.waitMessage('Processing...');
// ... long operation ...
Utils.waitMessageEnd();
\end{lstlisting}

\noindent
Wait messages stack: calling \texttt{waitMessage()} while one is already shown replaces the message text; calling \texttt{waitMessageEnd()} restores the previous message.

\begin{notebox}
\texttt{Server.call()} automatically shows a wait cursor during requests, so you only need \texttt{waitMessage()} for custom long operations or when you want a specific message.
\end{notebox}

